\section{Review methodology}\label{sec:methodology}
This section documents the procedures by which the review was conducted, following the evidence-based software engineering guidelines~\cite{KitchenhamCharters2007SLR}.
A pre-registered protocol is available at~\cite{kavadias2026qplSurveyProtocol}.
The section covers search protocol, screening and selection, study quality assessment, data extraction and synthesis and sensitivity analysis with threats to validity.
Results are reported in Section~\ref{sec:system_assessments}.

\subsection{Search protocol}\label{subsec:methodology:search_protocol}
Searches were conducted on 2026-04-27 across four bibliographic databases: ACM Digital Library, IEEE Xplore, arXiv and Google Scholar.
Official project and specification sites were consulted on a per-system basis throughout the assessment period.
A literature cutoff of 2026-04-08 was applied; evidence published after this date was not used for scoring.

Three verbatim search strings were applied consistently across all databases.
The primary string (S1) used broad seed terms targeting quantum programming languages and intermediate representations: \texttt{"quantum programming language" OR "quantum intermediate representation" OR OpenQASM OR "quantum MLIR" OR Silq OR Qiskit OR PennyLane}.
A second string (S2) targeted language-specific terms not reliably surfaced by domain vocabulary: \texttt{Silq quantum OR Quff quantum OR Qutes quantum OR Quipper quantum OR "Microsoft Q\#" OR "Q-MLIR"}.
A third string (S3) focused on abstraction-level framing: \texttt{"quantum programming" AND "abstraction" AND "language"}.
Full verbatim strings, per-database result counts and search URLs are published in Appendix~A of~\cite{kavadias2026qplSurveySupplementary}.

Pre-deduplication totals were 923 records from arXiv, 1,328 from ACM Digital Library, 728 from IEEE Xplore and 29,202 from Google Scholar, giving a combined total of 32,181 records.
The estimated unique record count after deduplication is 895--1,500; the Google Scholar corpus overlaps substantially with the other three databases.
A supplementary AI-assisted discovery process ran across the study period using a scheduled agent that scanned active project context and surfaced candidate papers on a fortnightly cadence.
This method does not produce reproducible verbatim queries and does not constitute a primary PRISMA search~\cite{page2021prisma}.
All candidates surfaced by this process were subsequently screened against the pre-specified criteria before inclusion.

\subsection{Screening and selection}\label{subsec:methodology:screening}
The review covers a pre-specified, bounded system set.
Candidate systems were assessed against four inclusion criteria.
A system is included if it meets at least two of the following:
\begin{description}
    \item[$C_1$] A formal specification or peer-reviewed primary paper exists.
    \item[$C_2$] Active releases or specification updates within 24 months prior to the cutoff.
    \item[$C_3$] Multivendor or multi-device backend support beyond a single lab or vendor.
    \item[$C_4$] Public documentation is sufficient to complete tasks defined in Section~\ref{subsec:evaluation_rubric:task_suite} without access to private knowledge.
\end{description}
Systems were excluded if their primary materials were marketing-only, if the project had been inactive for more than 36 months or if documentation was not enough to assess representative workflows.
The version rule stipulates that the latest stable major release or specification available by 2026-04-08 is the evaluation target; earlier versions are used only as labelled historical contrast.

Screening was conducted by a single reviewer.
This review covers a pre-defined system set rather than a large open literature corpus; full-text exclusion rationale is recorded for each system in Appendix~B of~\cite{kavadias2026qplSurveySupplementary}.
Two additional reviewers participate in independent scoring and calibration reconciliation.
Dual screening of the full inclusion corpus was not feasible given the review scope and is acknowledged as a limitation per PRISMA 2020 item~16b~\cite{page2021prisma}.

Ten systems were included in the primary review: OpenQASM~3.0, QIR, Q-MLIR, Qiskit, PennyLane, Q\#, Silq, Qmod, Quff and Qutes.
Quipper was excluded as a primary system, failing C\textsubscript{2} and C\textsubscript{3}; it is retained as a labelled historical contrast reference in Section~\ref{sec:system_descriptions}.

\subsection{Study quality assessment}\label{subsec:methodology:quality}
Evidence quality was assessed for each system prior to data extraction using a source-tier hierarchy adapted from the evidence-based software engineering guidelines~\cite{KitchenhamCharters2007SLR}.
Source tier~1 covers peer-reviewed journal or conference publications.
Source tier~2 covers preprints from established research groups with a peer-reviewed publication programme.
Source tier~3 covers official specifications, standards-body publications and maintained documentation.
No system in this review depends on source tier~4 grey literature as its primary source.

IR-tier systems, specifically QIR and Q-MLIR, use source tier~3 specification documentation as their primary and most authoritative source by design.
The field does not produce peer-reviewed primary papers for every IR or low-level framework; the absence of a standalone academic publication for QIR reflects the nature of the artefact rather than a gap in evidence retrieval.

Source-tier assignments, evidence base summaries and leave-one-out flags are given in Table~\ref{tab:quality_assessment}.
A \textit{High} confidence rating denotes a source tier~1 primary citation with two or more independent corroborating sources.
A \textit{Medium} confidence rating denotes a source tier~2 primary citation or a system with a single corroborating source.

\begin{table*}[htbp*]
    \centering
    \caption{Source quality assessment per system.
    LOO = leave-one-out sensitivity flag; see Subsection~\ref{subsec:methodology:sensitivity}.}
    \label{tab:quality_assessment}
    \begin{tabular}{p{1.5cm} p{5.5cm} c c c}
        \toprule
        \textbf{System} & \textbf{Primary citation(s)} & \textbf{Source tier} & \textbf{Confidence} & \textbf{LOO} \\
        \midrule
        OpenQASM   & \cite{cross2022openqasm3}                          & 1              & High   & ---        \\
        QIR        & \cite{qir_alliance_qir_spec_repo}                  & 3\dag          & Medium & ---        \\
        Q-MLIR     & \cite{McCaskey2021MLIRDialectQASM}                 & 1              & High   & ---        \\
        Qiskit     & \cite{aleksandrowicz2019qiskit}                    & 1              & High   & ---        \\
        PennyLane  & \cite{bergholm2018pennylane}                       & 1              & High   & ---        \\
        Q\#        & \cite{MicrosoftQSharpOverview}                     & 1, 3           & High   & ---        \\
        Silq       & \cite{bichsel2020silq}                             & 1              & High   & ---        \\
        Qmod       & \cite{vax2025qmodexpressivehighlevelquantum}       & 2              & Medium & ---        \\
        Quff       & \cite{wright2024quff}                              & 1              & Medium & \checkmark \\
        Qutes      & \cite{faro2025qutes,reversiblelifetime2026}        & 1              & High   & \checkmark \\
        \bottomrule
    \end{tabular}
    \smallskip

    \footnotesize\dag~Source tier~3 is the expected primary source for IR-tier artefacts; see text.
\end{table*}

\subsection{Data extraction and synthesis}\label{subsec:methodology:extraction}
For each system, the following data were extracted.
System identity, version, primary citation and documentation type were recorded alongside taxonomy category and typical workflow description.
Task completion evidence was captured for T\textsubscript{1}--T\textsubscript{5}, with axis scores A--E supported by one to three evidence quotes per axis citing the specific paper section or documentation anchor.
Quantum bias items with justification, confidence tier, scoring reviewer and date completed the extraction record.
Where evidence was absent or ambiguous, the more conservative score was assigned; scores inferred from examples marked as optional or advanced were not used.
The extraction dataset is published as Appendix~C and per-axis evidence quotes as Appendix~D of~\cite{kavadias2026qplSurveySupplementary}.

Calibration was conducted by two reviewers independently scoring three to five systems before proceeding to full independent scoring.
Disagreements greater than one level were discussed and resolved with a documented rationale; any amendments made post-calibration were applied prospectively.

Synthesis follows the structured narrative approach without meta-analysis~\cite{KitchenhamCharters2007SLR}.
Systems are grouped by the taxonomy tiers defined in Section~\ref{sec:taxonomy}.
Per-system outputs follow the aggregation rules defined in Section~\ref{sec:evaluation_rubric}: headline score, spectrum score (arithmetic mean of A--E with minimum--maximum range) and quantum bias count.

\subsection{Sensitivity analysis and threats to validity}\label{subsec:methodology:sensitivity}
A pre-specified leave-one-out sensitivity analysis was conducted on systems where a single document is the sole or pivotal source for one or more axis scores~\cite{kavadias2026qplSurveyProtocol}.
Two cases were identified.

\textit{Quff}: all axis scores depend on a single peer-reviewed source~\cite{wright2024quff}.
No additional published specification or paper existed at the review cutoff.
Removal of this source would leave no evidence base.
Scores are retained given the peer-reviewed venue; the single-source dependency is noted as a limitation.

\textit{Qutes axis~C and axis~D scores of~3}: both scores depend critically on one preprint~\cite{reversiblelifetime2026} for the reversible lifetime semantics mechanism.
The HPDC~2025 paper alone~\cite{faro2025qutes} would not sustain either score; removal would reduce both to score~2, shifting the headline from high-level to embedded framework.
The broader Qutes research programme has two peer-reviewed venue publications; the preprint status of this specific mechanism is noted as a limitation rather than a basis for score reduction.

The principal threats to validity are as follows.
\begin{itemize}
    \item \textit{Scoring subjectivity}: is mitigated by dual independent scoring with calibration, required evidence quotes and predefined decision rules for ambiguous cases.
    \item \textit{Rubric dependence on T\textsubscript{1}--T\textsubscript{5}}: the task suite was pre-specified in the protocol~\cite{kavadias2026qplSurveyProtocol}; any post-calibration amendments are documented and applied prospectively only.
    \item \textit{Maturity differences across ecosystems}: source-tier labelling and per-claim confidence ratings mitigate misattribution.
    \item \textit{Version churn}: the versioning rule, latest stable major release by 2026-04-08, is applied consistently.
    \item \textit{Publication bias}: grey literature, including official specifications, standards-body publications and maintainer documentation, is included by design to avoid systematic under-representation of IR-tier and industry-produced frameworks~\cite{KitchenhamCharters2007SLR,page2021prisma}.
\end{itemize}
